\documentclass{article}
\usepackage{geometry}
\geometry{a4paper, margin=1in}
\usepackage{graphicx}
\usepackage{hyperref}

\title{Practical Machine Learning and Deep Learning\\D1.3 Progress Submission\\
Team "Magic Photo"}
\author{Maksim Ilin \href{mailto:m.ilin@innopolis.university}{m.ilin@innopolis.university},\\ Rail Sabirov \href{mailto:rai.sabirov@innopolis.university}{rai.sabirov@innopolis.university},\\ Ivan Ilyichev \href{mailto:i.ilyichev@innopolis.university}{i.ilyichev@innopolis.university}}

\date{\today}

\begin{document}

\maketitle

\section{Project Topic}
\par This project objective is to enhance existing solutions for image restoration from different corruptions and coloring of grayscale images using machine learning and deep learning techniques. The project will focus on developing a user-friendly application that allows users to upload images, select the processing option they want to apply, and receive the processed image as output. Additionally, we plan to collect users' feedback for future product improvements and keep track of software lifecycle. The application will leverage pre-trained models and fine-tune them. 

\section{Link to the Repository}
\par The link: \href{https://github.com/ivaniinno/Image-Restoration-and-Coloring/tree/dev}{Main repository}. For now we store code in \textbf{dev} branch.
\begin{itemize}
    \item \textbf{eda} folder contains all the steps for exploring and preparing the chosen datasets
    \item \textbf{data\_augmentation\_methods} folder contains image manipulation (preparation) script for training
    \item \textbf{docs} folder contains project documentation
    \item \textbf{models\_testing} folder contains investigation of baseline model implementations
    \item \textbf{models} folder contains 2 models that we currently finetune (pix2pix colorizer and real-esrgan for image restoration), for usage simplicity (to be described below) two standalone repositories are available for them: \href{https://github.com/RailSAB/coloringModel}{Colorization repository} and \href{https://github.com/RailSAB/real-esrganmodel}{Real-ERSGAN repository}
\end{itemize} 

\section{Dataset Splitting}
\par Since we have in total more than 100.000 images, we decided to split the dataset into 5 equal portions and train models for the same number of epochs on each partition to meet the resource limitations that are present on kaggle.
\par The list of splits is as follows:

\begin{enumerate}
    \item \href{https://www.kaggle.com/datasets/railsabirov/trainlid-split-1fifth}{1st partition}
    \item \href{https://www.kaggle.com/datasets/railsabirov/TrainLID-Split-1fifthV2/data}{2nd partition}
    \item \href{https://www.kaggle.com/datasets/railsabirov/TrainLID-Split-1fifthV3/data}{3rd partition}
    \item \href{https://www.kaggle.com/datasets/railsabirov/TrainLID-Split-1fifthV4/data}{4th partition}
    \item \href{https://www.kaggle.com/datasets/railsabirov/TrainLID-Split-1fifthV5/data}{5th partition}
\end{enumerate}

\section{Model Finetuning}
\par We have been finetuning both models with configured scripts on kaggle: \href{https://www.kaggle.com/code/railsabirov/colorizationmodelfinetune}{Colorization notebook} and \href{https://www.kaggle.com/code/railsabirov/real-esrgan-training}{Real-ESRGAN notebook}.
\par  Since we have encountered configuration issues related to python package versions, we created two standalone repositories for simpler usage of baseline models on kaggle: \href{https://github.com/RailSAB/coloringModel}{Colorization repository} and \href{https://github.com/RailSAB/real-esrganmodel}{Real-ERSGAN repository}. Additionally, we were bounded by kaggle resources (e.g., batch size of 4 is the highest available for restoration model, therefore time for 1 epoch on 20\% size of the dataset varies between 1,5 and 2 hours with GPU accelerator on).
\subsection{Intermediate Results}
\par The weights obtained during training are available as follows:

\begin{itemize}
    \item For colorization model: \href{https://www.kaggle.com/models/railsabirov/coloring}{Colorization Weights} (4 epochs/checkpoints are already stored)
    \item For restoration model: \href{https://www.kaggle.com/datasets/ivanilichev/checkpoints-for-readl-ersgan}{Restoration Weights} (11 epochs/checkpoints are already stored)
\end{itemize}


\section{Future Work}
\par Before the final submissions we have plans as follows:

\begin{itemize}
    \item Continue finetuning colorization and restoration models to achieve at least 4-5 epochs on each 20\% dataset split (resulting in 20-25 epochs in total).
    \item Compare 0th and last epochs for each model with key findings and discussion as a result.
    \item Create a service for interaction with the finetuned models.
\end{itemize}

\end{document}